\section{Algorithmic Implications}

This chapter records the algorithmic reading of the local-complexity
parameters used throughout the manuscript.  The claims here are
interpretive: they organize the consequences of the structural bounds
rather than introducing a new complexity-theoretic separation.

\subsection{Bounded refinement as bounded parallel depth}

Small $\chi_R$ indicates that the refinement process can be scheduled
with uniformly bounded local conflict.  In this regime, independent
local updates can be grouped into a bounded number of parallel rounds.

\begin{lemma}[Bounded scheduling interpretation]
If a family has uniformly bounded refinement chromatic parameter
$\chi_R$, then the associated local refinement schedule has uniformly
bounded round complexity, subject to the capacity assumptions already
fixed in the preceding chapters.
\end{lemma}

\subsection{Growth as a sequentiality obstruction}

Large $\chi_R$ records the opposite phenomenon: local conflicts cannot
be eliminated by a uniformly bounded schedule.  The obstruction is not
a lower bound for every algorithm, but a lower bound for the refinement
model under the stated locality and capacity rules.

\begin{definition}[Local-complexity obstruction]
A family exhibits a local-complexity obstruction when its refinement
conflict parameter grows without a uniform bound along the family.
\end{definition}

\subsection{Scope}

The chapter does not claim an unconditional complexity-class
separation.  It only records how the manuscript's refinement parameters
translate into bounded-depth and non-bounded-depth scheduling language.
